\begin{abstract}
This dissertation explores new perspectives on quantum technologies, advancing both quantum computation and quantum sensing through theoretical modeling and numerical simulation. The work addresses three challenges in the quantum technology landscape: hardware optimization for elevated-temperature operation, hybrid quantum-classical machine learning architectures, and quantum-enhanced single-photon detection.

The first contribution presents a comprehensive theoretical analysis of asymmetric Superconducting Quantum Interference Device (SQUID) transmon qubits based on NbN/AlN/NbN Josephson junctions, designed for operation at 4~K rather than conventional millikelvin temperatures. Through detailed electromagnetic simulations and decoherence modeling, we demonstrate that flux-tunable sweet spots enable dynamic frequency control while maintaining coherence times exceeding 50~$\mu$s. This architecture significantly reduces cryogenic complexity and energy overhead, paving the way for scalable, cryocooler-based quantum processors.

The second contribution develops and benchmarks hybrid quantum-classical convolutional neural networks (Q-CNNs) for satellite image classification using the EuroSAT dataset. Our simulations seem to reveal that quantum variational circuits accelerate training convergence by approximately 3$\times$ compared to classical CNNs, achieving $>$90\% accuracy within 3--5 epochs versus 15--20 epochs classically. This could lead to demonstrating that accelerated training efficiency, rather than asymptotic accuracy, may represent a primary quantum advantage in near-term applications. Hardware validation on IQM's 54-qubit processor will hopefully confirm functional correctness.

The third contribution models Rydberg atom-based avalanche detectors for single-photon sensitivity. Numerical simulations across 0--5~T magnetic fields predict amplification factors exceeding 3, and dark count rates below 0.01~Hz at 4~K, with robust detection time stability. This work establishes theoretical performance bounds for quantum-enhanced weak-signal detection relevant to metrology and fundamental physics searches.

Collectively, these contributions demonstrate that simulation-driven research enables systematic exploration of quantum parameter spaces, identifies practical regimes of quantum advantage, and guides experimental implementations toward deployable quantum technologies. The dissertation emphasizes cross-cutting insights, including hardware tunability as a design principle, resource efficiency as a near-term quantum value proposition, and the convergence of sensing and computation in hybrid quantum systems operating across temperature scales and physical platforms.
\end{abstract}
